% Figure: the mean-variance efficient frontier traced by a Lagrange
% problem. Minimising portfolio variance at each target return, over
% weights that sum to one, sweeps out a parabola in the
% (variance, return) plane; its left edge is the minimum-variance
% portfolio and its upper branch is the efficient frontier. The
% Lagrange multiplier on the return constraint is the slope of the
% frontier in the (variance, return) plane.
\begin{figure}[htbp]
\centering
\begin{tikzpicture}[scale=1.0,
  eff/.style={engblue, very thick},
  ineff/.style={engblue, very thick, dashed},
  guide/.style={enggray, dashed},
  pt/.style={circle, fill=engblack, inner sep=1.5pt},
]
  \begin{axis}[
    width=9.6cm, height=6.6cm,
    xlabel={portfolio variance $\sigma_{p}^{2}$},
    ylabel={expected return $\mu_{p}$},
    xmin=0, xmax=0.09, ymin=0.02, ymax=0.16,
    axis lines=left,
    xtick={0,0.02,0.04,0.06,0.08},
    ytick={0.04,0.08,0.12,0.16},
    tick label style={font=\small},
    label style={font=\small},
    clip=false,
  ]
    % Frontier: variance is quadratic in return, sigma^2 = a(mu - mu0)^2 + s0.
    % Take mu0 = 0.08 (min-variance return), s0 = 0.01 (min variance),
    % a = 4. Upper (efficient) branch mu >= 0.08, lower (inefficient) branch.
    % Efficient branch (solid): mu from 0.08 up to 0.155.
    \addplot[eff, domain=0.08:0.155, samples=60]
      ({4*(x-0.08)^2 + 0.01}, {x});
    % Inefficient branch (dashed): mu from 0.03 to 0.08.
    \addplot[ineff, domain=0.03:0.08, samples=40]
      ({4*(x-0.08)^2 + 0.01}, {x});
    % Minimum-variance portfolio marker.
    \node[pt] at (axis cs:0.01,0.08) {};
    \node[anchor=west, font=\small] at (axis cs:0.012,0.075)
      {min-variance};
    % A target-return operating point on the efficient branch.
    % mu = 0.12: sigma^2 = 4*(0.04)^2 + 0.01 = 4*0.0016+0.01 = 0.0164.
    \node[pt] at (axis cs:0.0164,0.12) {};
    \draw[guide] (axis cs:0.0164,0.02) -- (axis cs:0.0164,0.12);
    \node[engblue, anchor=south west, font=\small] at (axis cs:0.045,0.135)
      {efficient frontier};
    \node[enggray, anchor=north west, font=\small] at (axis cs:0.03,0.06)
      {inefficient};
  \end{axis}
\end{tikzpicture}
\caption[The mean-variance efficient frontier as a swept Lagrange
solution.]{Minimising portfolio variance at each fixed target return, over
asset weights constrained to sum to one, is a Lagrange problem whose
solution, swept across all target returns, traces a parabola in the
variance-return plane. The leftmost point is the minimum-variance
portfolio; the solid upper branch is the efficient frontier, the set of
portfolios with the least variance for their return, and the dashed lower
branch is dominated, since each of its portfolios shares a variance with a
higher-return portfolio directly above it. The Lagrange multiplier on the
return constraint is the marginal variance the investor pays for one more
unit of return, read off as the local steepness of the frontier.}
\label{fig:vol02:ch11:mean-variance-frontier}
\end{figure}
